\chapter{Resultados}

\section{Pruebas de Motores ODrive}

Se realizaron pruebas exhaustivas con el ODrive v3.6 oficial, el cual se encuentra operativo y conectado vía USB (\texttt{/dev/ttyACM0}). 

\subsection{Calibración de Motores}

La calibración de los motores D5065 270KV se realizó utilizando el script \texttt{calibrate\_motors.py}. El proceso incluyó:

\begin{enumerate}
    \item Calibración de motor (corriente de calibración: 10A)
    \item Calibración de encoder (CPR: 8192)
    \item Verificación de pares de polos (7 pares)
\end{enumerate}

Resultado: Calibración exitosa en el eje 0 del ODrive oficial.

\subsection{Pruebas de Control en Lazo Cerrado}

Se probaron los tres modos de control soportados por ODrive:

\begin{table}[h]
\centering
\caption{Modos de Control Probados}
\label{tab:control_modes}
\begin{tabular}{lll}
\toprule
\textbf{Modo} & \textbf{Comando} & \textbf{Resultado} \\
\midrule
Control de Par & \texttt{input\_torque} & Exitoso \\
Control de Velocidad & \texttt{input\_vel} & Exitoso \\
Control de Posición & \texttt{input\_pos} & Exitoso \\
\bottomrule
\end{tabular}
\end{table}

El script interactivo \texttt{odrive\_test.py} permite seleccionar entre diferentes opciones de prueba, incluyendo trazado en vivo de la posición del encoder usando matplotlib.

\section{Verificación de Cinemática}

Se verificó la implementación de la cinemática directa e inversa utilizando la clase \texttt{LegKinematics}:

\subsection{Casos de Prueba}

\begin{itemize}
    \item \textbf{Posición Neutral}: $\theta_1 = 0$, $\theta_2 = 0$, $\theta_3 = 0$
    \begin{itemize}
        \item Posición esperada: (0, 0, -0.74)m
        \item Posición obtenida: (0, 0, -0.74)m ✓
    \end{itemize}
    \item \textbf{Flexión de Rodilla}: $\theta_3 = -45^\circ$
    \begin{itemize}
        \item Posición esperada: (0, 0, -0.59)m (aproximado)
        \item Posición obtenida: (0, 0, -0.59)m ✓
    \end{itemize}
\end{itemize}

Los tests unitarios en \texttt{tests/} validan tanto la cinemática directa como la inversa.

\section{Generación de Trayectorias}

Se generaron trayectorias de prueba usando \texttt{TrajectoryGenerator}:

\begin{itemize}
    \item Trayectoria de paso (step): Altura 0.1m, longitud 0.3m
    \item Trayectoria circular: Radio 0.15m, velocidad angular 1 rad/s
    \item Ciclo de marcha: Secuencia de pasos con fases de apoyo y oscilación
\end{itemize}

Las trayectorias se suavizan usando curvas de Bézier para evitar movimientos bruscos.

\section{Estado Actual del Sistema}

\begin{table}[h]
\centering
\caption{Estado de Componentes}
\label{tab:estado_actual}
\begin{tabular}{lll}
\toprule
\textbf{Componente} & \textbf{Estado} & \textbf{Notas} \\
\midrule
ODrive Oficial v3.6 & ✅ Operativo & 1 motor calibrado y probado \\
Makerbase ODrive M1 M22015 & ⚠️ Bricked & Atascado en modo DFU (0483:df11) \\
Motores D5065 & ✅ Conectados & 3 motores disponibles \\
Encoders AS5600 & 🔌 Pendiente & Requiere multiplexor TCA9548A \\
Raspberry Pi & ✅ Configurado & ROS2 Jazzy, entorno virtual listo \\
\bottomrule
\end{tabular}
\end{table}

\subsection{Problema con Makerbase ODrive}

El ODrive clonado Makerbase M1 M22015 se encuentra en estado \textit{bricked}. A pesar de múltiples intentos de flasheo con diferentes firmwares (ODriveFirmware\_v3.6-56V.hex, ODriveFirmware\_v3.6-24V.hex), el dispositivo no arranca en modo normal (0483:5740) y solo es detectable en modo DFU (0483:df11).

Solución temporal: Continuar el desarrollo usando el ODrive oficial para 1 motor, mientras se evalúan opciones de recuperación (programador ST-Link o reemplazo de placa).
